\chapter{Literature Review}
\label{ch:lit_review}

% ---------------------------------------------------------------------
% Write this chapter so it can be lifted into a standalone
% "related work" / background section of a journal or conference
% paper later with minimal changes:
%   - avoid thesis-only meta-references ("as discussed in Chapter 4")
%   - keep subsections self-contained and topic-based, not
%     chronological narration
%   - every claim traces back to a row in your concept-matrix
%     (bib/concept_matrix.csv) and a \cite{} key in references.bib
% ---------------------------------------------------------------------

\section{Scope and search strategy}
% Databases searched, search terms, inclusion/exclusion criteria,
% snowballing approach (see Masters_Prep_Plan.md, Section 17).

\section{Background}
% Domain background a reader needs before the related work
% (e.g. histology, image registration fundamentals).

\section{Related Work}
% Organise by sub-theme, not by paper. Each subsection lives in its own
% file under chapters/lit_review/ so they can be drafted, reviewed and
% (later) lifted into a standalone paper independently.
\subsection{Classical (non-learning) registration methods}
\label{subsec:classical_registration}

% Classical/optimisation-based registration: transformation models
% (rigid, affine, B-spline/FFD, demons), similarity metrics (SSD, NCC,
% mutual information), and regularisation. See Sotiras et al. (2013)
% and Pluim et al. (2003) for the two key surveys.


\subsection{Deep-learning-based registration methods}
\label{subsec:dl_registration}

% Learning-based registration: supervised vs. unsupervised, VoxelMorph
% and successors, transformer-based methods, histology-specific
% adaptations (e.g. ANHIR/ACROBAT challenge methods).


\subsection{Evaluation methodologies}
\label{subsec:evaluation_methodologies}

% How registration quality is measured: target registration error
% (TRE), Dice overlap, Jacobian determinant (plausibility), statistical
% comparison methods (Wilcoxon signed-rank, median/IQR reporting).


\section{Research Gap}
% The blank/weak cells in your concept matrix. This section is the
% hinge between the lit review and Chapter~\ref{ch:methodology}.
